\section{Experiments}
\label{sec:experiments}

\subsection{Setup}

We use the processed OWA version of ProofWriter. After flattening each theory into question-level instances, the \texttt{depth-3ext-NatLang} split contains 112,062 training questions, 17,312 development questions, and 34,378 test questions. The \texttt{depth-5} split contains 36,902 training questions, 10,190 development questions, and 20,030 test questions. For transfer, we also evaluate on \texttt{NatLang/test}, \texttt{depth-3/test}, and \texttt{birds-electricity/test}.

Our study has two parts. First, we run a seeded 7B suite at a fixed 4,096-question training budget. The training seed changes the sampled training questions, while the evaluation subset is held fixed with data seed 0. This lets us measure model variance without changing the test slice. Second, we scale the same comparison to larger training budgets and evaluate selected checkpoints on full test splits for transfer analysis.

We compare three main supervision formats:
\begin{itemize}
\item \textbf{answer-only}: predict only the label,
\item \textbf{proof-only}: emit structured reasoning without explicit abstention-witness supervision,
\item \textbf{\model{}}: emit proof, refutation, or abstention witnesses.
\end{itemize}
We also run three single-seed ablations at the same 4k budget:
\begin{itemize}
\item \textbf{ProCo-chain}: keep abstention failure tags but replace detailed witnesses with a generic note,
\item \textbf{ProCo-witness}: keep detailed missing-support witnesses but remove explicit failure tags,
\item \textbf{ProCo-no-refute}: keep the \model{} target for true and unknown queries, but remove explicit refutation evidence on false cases.
\end{itemize}

All runs use QLoRA with NF4 4-bit loading, LoRA rank 16, LoRA alpha 32, dropout 0.05, maximum sequence length 512, learning rate \(2\times 10^{-4}\), and one training epoch. We report label accuracy, macro-F1, Unknown-class F1, evidence faithfulness, and joint accuracy. Joint accuracy requires both a correct label and faithful evidence.

\subsection{0.5B Pilot: Isolating the Supervision Effect}

Table~\ref{tab:pilot-results} summarizes the initial 0.5B pilot. At this scale, answer-only is still the strongest raw classifier. Proof-only gives up raw accuracy and produces almost no faithful evidence. \model{} recovers some of the lost accuracy and greatly improves faithfulness relative to proof-only, but it does not close the raw-accuracy gap to answer-only.

\begin{table}[t]
\centering
\small
\begin{tabular}{llccccc}
\toprule
Domain & Variant & Acc. & Macro-F1 & Unknown F1 & Faithful & Joint \\
\midrule
ID & answer-only & 74.7 & 75.1 & 71.1 & 0.0 & 0.0 \\
ID & proof-only & 70.7 & 71.3 & 62.5 & 5.3 & 5.2 \\
ID & \model{} & 71.3 & 71.8 & 66.3 & 29.6 & 29.4 \\
\midrule
OOD & answer-only & 64.2 & 64.7 & 61.7 & 0.0 & 0.0 \\
OOD & proof-only & 63.0 & 63.9 & 54.9 & 0.8 & 0.8 \\
OOD & \model{} & 61.8 & 62.7 & 54.4 & 20.2 & 20.2 \\
\bottomrule
\end{tabular}
\caption{0.5B pilot results on ProofWriter. At small scale, \model{} improves verifier-backed reasoning relative to proof-only supervision but still trails answer-only in raw accuracy.}
\label{tab:pilot-results}
\end{table}

We use this pilot as motivation only. The paper's main claim comes from the 7B seeded and scaling results below.

\subsection{7B Seeded Results on a Fixed Test Subset}

Table~\ref{tab:seeded-results} is the main robustness result. At the 4k budget, answer-only still has the highest raw accuracy, but \model{} gives a much larger gain in verifier-backed reasoning. In-domain, joint accuracy rises from \(39.9 \pm 0.6\) to \(75.0 \pm 0.7\), while raw accuracy changes only from \(87.9 \pm 0.2\) to \(87.4 \pm 0.3\). On depth-OOD, joint accuracy rises from \(26.5 \pm 1.1\) to \(55.0 \pm 0.7\), with nearly unchanged raw accuracy (\(76.1 \pm 0.7\) for proof-only and \(76.4 \pm 0.6\) for \model{}).

\begin{table}[t]
\centering
\small
\begin{tabular}{llccccc}
\toprule
Domain & Variant & Runs & Acc. & Unknown F1 & Faithful & Joint \\
\midrule
ID & answer-only & 3 & 88.9 $\pm$ 0.3 & 87.7 $\pm$ 0.4 & 0.0 $\pm$ 0.0 & 0.0 $\pm$ 0.0 \\
ID & proof-only & 3 & 87.9 $\pm$ 0.2 & 86.7 $\pm$ 0.3 & 40.0 $\pm$ 0.6 & 39.9 $\pm$ 0.6 \\
ID & \model{} & 3 & 87.4 $\pm$ 0.3 & 86.0 $\pm$ 0.6 & 75.1 $\pm$ 0.7 & 75.0 $\pm$ 0.7 \\
\midrule
Depth-OOD & answer-only & 3 & 80.0 $\pm$ 1.0 & 76.3 $\pm$ 0.6 & 0.0 $\pm$ 0.0 & 0.0 $\pm$ 0.0 \\
Depth-OOD & proof-only & 3 & 76.1 $\pm$ 0.7 & 73.0 $\pm$ 0.5 & 26.6 $\pm$ 1.1 & 26.5 $\pm$ 1.1 \\
Depth-OOD & \model{} & 3 & 76.4 $\pm$ 0.6 & 73.3 $\pm$ 1.0 & 55.1 $\pm$ 0.7 & 55.0 $\pm$ 0.7 \\
\bottomrule
\end{tabular}
\caption{Seeded 7B results on a fixed sampled test subset.}
\label{tab:seeded-results}
\end{table}

Relative to answer-only, the raw-accuracy gap is 1.5 points in-domain and 3.6 points on depth-OOD, but answer-only does not provide evidence at all. At this training budget, the right reading is therefore a small raw-accuracy loss in exchange for a large and stable gain in checkable reasoning.

\subsection{Additional Backbone Check}

To test whether the effect is tied to Qwen2.5-7B, we also run the same 4k
training comparison with Mistral-7B-Instruct-v0.3 for seed 0. Table~\ref{tab:backbone-check}
shows the same pattern, and in this run the raw-accuracy tradeoff is smaller:
\model{} matches answer-only in-domain within 0.1 points and exceeds it on
depth-OOD, while joint accuracy remains far above proof-only. This is not yet a
multi-seed backbone study, so we treat it as a robustness check rather than a
headline result.

\begin{table}[t]
\centering
\small
\begin{tabular}{lllcccc}
\toprule
Model & Domain & Variant & Runs & Acc. & Faithful & Joint \\
\midrule
Qwen2.5-7B & ID & answer-only & 3 & 88.9 $\pm$ 0.3 & 0.0 $\pm$ 0.0 & 0.0 $\pm$ 0.0 \\
Qwen2.5-7B & ID & proof-only & 3 & 87.9 $\pm$ 0.2 & 40.0 $\pm$ 0.6 & 39.9 $\pm$ 0.6 \\
Qwen2.5-7B & ID & ProCo & 3 & 87.4 $\pm$ 0.3 & 75.1 $\pm$ 0.7 & 75.0 $\pm$ 0.7 \\
\cmidrule(lr){2-7}
Qwen2.5-7B & Depth-OOD & answer-only & 3 & 80.0 $\pm$ 1.0 & 0.0 $\pm$ 0.0 & 0.0 $\pm$ 0.0 \\
Qwen2.5-7B & Depth-OOD & proof-only & 3 & 76.1 $\pm$ 0.7 & 26.6 $\pm$ 1.1 & 26.5 $\pm$ 1.1 \\
Qwen2.5-7B & Depth-OOD & ProCo & 3 & 76.4 $\pm$ 0.6 & 55.1 $\pm$ 0.7 & 55.0 $\pm$ 0.7 \\
\midrule
Mistral-7B & ID & answer-only & 1 & 87.7 & 0.0 & 0.0 \\
Mistral-7B & ID & proof-only & 1 & 86.4 & 42.0 & 41.9 \\
Mistral-7B & ID & ProCo & 1 & 87.6 & 79.2 & 79.0 \\
\cmidrule(lr){2-7}
Mistral-7B & Depth-OOD & answer-only & 1 & 77.8 & 0.0 & 0.0 \\
Mistral-7B & Depth-OOD & proof-only & 1 & 77.3 & 31.6 & 31.5 \\
Mistral-7B & Depth-OOD & ProCo & 1 & 78.6 & 61.6 & 61.4 \\
\bottomrule
\end{tabular}
\caption{Additional backbone check at 4k training examples. Qwen2.5-7B rows use
three seeds; Mistral-7B rows are a single seed.}
\label{tab:backbone-check}
\end{table}

\subsection{Where the Gain Comes From: Abstention Evidence}

Table~\ref{tab:abstention-evidence} looks only at gold-unknown examples and
separates three outcomes: a faithful abstention, a generic or unfaithful
\texttt{Unknown}, and an over-commitment to \texttt{True} or \texttt{False}. The
key point is that \model{} does not simply abstain more often. In-domain,
proof-only produces \(1544.0 \pm 30.4\) generic unknown outputs and zero
faithful unknown explanations. \model{} reduces generic unknown outputs to
\(81.0 \pm 36.9\) and produces \(1444.3 \pm 7.4\) faithful missing-support
witnesses. On depth-OOD, proof-only again produces zero faithful unknown
explanations, while \model{} produces \(1169.0 \pm 29.0\).

\begin{table}[t]
\centering
\small
\begin{tabular}{llcccc}
\toprule
Domain & Variant & Faithful & Generic & Over-Commit & Faithful Rate \\
\midrule
ID & proof-only & 0.0 $\pm$ 0.0 & 1544.0 $\pm$ 30.4 & 316.0 $\pm$ 30.4 & 0.0 $\pm$ 0.0 \\
ID & \model{} & 1444.3 $\pm$ 7.4 & 81.0 $\pm$ 36.9 & 334.7 $\pm$ 39.6 & 77.7 $\pm$ 0.4 \\
\midrule
Depth-OOD & proof-only & 0.0 $\pm$ 0.0 & 1278.0 $\pm$ 56.8 & 426.0 $\pm$ 56.8 & 0.0 $\pm$ 0.0 \\
Depth-OOD & \model{} & 1169.0 $\pm$ 29.0 & 117.0 $\pm$ 43.7 & 418.0 $\pm$ 72.7 & 68.6 $\pm$ 1.7 \\
\bottomrule
\end{tabular}
\caption{Gold-unknown behavior on the fixed 7B test subset. Faithful means the
model predicts \texttt{Unknown} and emits a verifier-checkable missing-support
witness. Generic means the model predicts \texttt{Unknown} but the witness is not
faithful. Over-commit means the model predicts \texttt{True} or \texttt{False}.}
\label{tab:abstention-evidence}
\end{table}

The gain therefore comes from replacing generic abstentions with verifiable
missing-support witnesses, not from changing the abstention rate alone.

\subsection{Support-Deletion Mutation Stress Test}

To test whether missing-support witnesses survive a controlled intervention, we
construct a support-deletion set from originally provable test questions. For
each source question with a \texttt{True} or \texttt{False} answer, we delete one
fact or rule from its gold derivation chain. We keep the mutant only when the
originally supported literal is no longer derivable and the opposite literal is
also not derivable. Each retained mutant is therefore a gold-\texttt{Unknown}
example with a known deleted support and a verifier-generated abstention
witness.

\begin{table}[t]
\centering
\small
\begin{tabular}{lcccc}
\toprule
Variant & Runs & Pred. Unknown & Faithful Unknown & Joint \\
\midrule
answer-only & 1 & 2110.0 & 0.0 & 0.0 \\
proof-only & 1 & 2294.0 & 0.0 & 0.0 \\
ProCo & 1 & 2127.0 & 1999.0 & 50.0 \\
\bottomrule
\end{tabular}
\caption{Support-deletion mutation stress test on 4,000 Qwen2.5-7B examples.
All examples are gold-\texttt{Unknown} after deleting one necessary fact or rule
from an originally provable source question. Joint is equivalent to faithful
unknown rate on this set.}
\label{tab:support-deletion}
\end{table}

Table~\ref{tab:support-deletion} shows that the mutation is not solved by
generic conservatism. Answer-only predicts \texttt{Unknown} on 2,110 of 4,000
mutants. Proof-only predicts \texttt{Unknown} more often, 2,294 of 4,000, but
still produces no faithful missing-support witnesses. \model{} predicts
\texttt{Unknown} on a similar number of mutants, 2,127 of 4,000, yet 1,999 of
those predictions contain verifier-checkable missing-support evidence. The
controlled deletion setting therefore supports the same mechanism as the
natural gold-unknown slice: the improvement is evidence quality, not a higher
abstention rate.

\subsection{Ablating the Abstention Target}

Because the biggest behavioral gap appears on unsupported queries, Table~\ref{tab:ablation-results} ablates the abstention target itself. Both partial variants recover most of the gain over proof-only. In-domain joint accuracy reaches 74.4 for ProCo-chain and 74.0 for ProCo-witness, close to 75.2 for full \model{} and far above 40.2 for proof-only. On depth-OOD, the corresponding numbers are 55.6, 53.4, 55.1, and 26.7.

\begin{table}[t]
\centering
\small
\begin{tabular}{llcccc}
\toprule
Domain & Variant & Runs & Acc. & Verifier-accepted & Joint \\
\midrule
ID & proof-only & 1 & 86.4 & 42.0 & 41.9 \\
ID & proof-only & 3 & 87.9 $\pm$ 0.2 & 40.0 $\pm$ 0.6 & 39.9 $\pm$ 0.6 \\
ID & \model{}-chain & 1 & 87.0 & 74.5 & 74.4 \\
ID & \model{}-witness & 1 & 87.9 & 74.1 & 74.0 \\
ID & \model{} & 1 & 87.6 & 79.2 & 79.0 \\
ID & \model{} & 3 & 87.4 $\pm$ 0.3 & 75.1 $\pm$ 0.7 & 75.0 $\pm$ 0.7 \\
\midrule
Depth-OOD & proof-only & 1 & 77.3 & 31.6 & 31.5 \\
Depth-OOD & proof-only & 3 & 76.1 $\pm$ 0.7 & 26.6 $\pm$ 1.1 & 26.5 $\pm$ 1.1 \\
Depth-OOD & \model{}-chain & 1 & 77.8 & 55.6 & 55.6 \\
Depth-OOD & \model{}-witness & 1 & 78.2 & 53.5 & 53.4 \\
Depth-OOD & \model{} & 1 & 78.6 & 61.6 & 61.4 \\
Depth-OOD & \model{} & 3 & 76.4 $\pm$ 0.6 & 55.1 $\pm$ 0.7 & 55.0 $\pm$ 0.7 \\
\bottomrule
\end{tabular}
\caption{Ablations on the abstention target.}
\label{tab:ablation-results}
\end{table}

The full target is still best or tied across the two domains. This suggests that failure tags and witness text help in different ways, and that the best results come from keeping both.

\subsection{Ablating Explicit Refutation}

To separate abstention supervision from explicit refutation evidence, we run a second single-seed 4k ablation. ProCo-no-refute keeps the \model{} target for true and unknown queries, but on false queries it predicts \texttt{LABEL: False} without a refutation chain. By construction, this variant cannot receive faithfulness credit on false examples. The question is therefore not whether it can match full \model{}, but how much of the gain survives when proof and abstention evidence are kept and refutation evidence is removed.

\begin{table}[t]
\centering
\small
\begin{tabular}{llccccc}
\toprule
Domain & Variant & Acc. & False F1 & Unknown F1 & Faithful & Joint \\
\midrule
ID & proof-only & 87.6 & 87.3 & 86.4 & 40.4 & 40.2 \\
ID & ProCo-no-refute & 86.3 & 85.7 & 84.6 & 53.1 & 53.1 \\
ID & ProCo & 87.1 & 86.8 & 85.5 & 75.3 & 75.2 \\
\midrule
Depth-OOD & proof-only & 77.0 & 77.5 & 73.3 & 26.9 & 26.7 \\
Depth-OOD & ProCo-no-refute & 76.2 & 76.9 & 72.0 & 39.9 & 39.8 \\
Depth-OOD & ProCo & 77.2 & 77.6 & 73.9 & 55.2 & 55.1 \\
\bottomrule
\end{tabular}
\caption{Single-seed 4k ablation of explicit refutation evidence. ProCo-no-refute keeps the \model{} target on true and unknown queries but removes refutation chains on false queries.}
\label{tab:refute-ablation-results}
\end{table}

Table~\ref{tab:refute-ablation-results} gives a mixed answer. On the one hand, ProCo-no-refute clearly improves over proof-only: joint accuracy rises from 40.2 to 53.1 in-domain and from 26.7 to 39.8 on depth-OOD, and faithful unknown explanations remain high at 1,378 in-domain and 1,093 on depth-OOD. On the other hand, it falls well short of full \model{}, which reaches 75.2 and 55.1 joint accuracy on the same slices. False F1 changes only modestly (85.7 vs.\ 86.8 in-domain; 76.9 vs.\ 77.6 on depth-OOD), but faithfulness drops sharply because false predictions no longer come with verifiable refutation traces. The cleanest reading is that abstention supervision explains much of the improvement on unsupported queries, while full verifier-backed performance still benefits substantially from explicit refutation evidence.

\subsection{Scaling and Transfer}

Tables~\ref{tab:scaling-results} and~\ref{tab:transfer-results} show what happens when the same comparison is run at larger training budgets and on additional test sets. Two trends are clear. First, faithfulness and joint accuracy rise steadily with more data. In-domain joint accuracy climbs from \(75.0\) at 4k to \(88.4\) at 16k, \(93.4\) at 32k, and \(96.8\) at full train. On depth-OOD, it climbs from \(55.0\) to \(71.2\), \(77.6\), and \(86.2\). Second, the raw-accuracy gap to answer-only shrinks with scale. In-domain, the gap narrows from 1.5 points at 4k to 0.3 points at full train. On the fixed depth-5 subset, \model{} slightly surpasses answer-only at full train, \(93.8\) versus \(93.1\).

\begin{table}[t]
\centering
\small
\begin{tabular}{rllccc}
\toprule
Train & Domain & Variant & Acc. & Faithful & Joint \\
\midrule
4096 & ID & answer-only & 88.9 $\pm$ 0.3 & 0.0 $\pm$ 0.0 & 0.0 $\pm$ 0.0 \\
4096 & ID & proof-only & 87.9 $\pm$ 0.2 & 40.0 $\pm$ 0.6 & 39.9 $\pm$ 0.6 \\
4096 & ID & ProCo & 87.4 $\pm$ 0.3 & 75.1 $\pm$ 0.7 & 75.0 $\pm$ 0.7 \\
\cmidrule(lr){2-6}
4096 & Depth-OOD & answer-only & 80.0 $\pm$ 1.0 & 0.0 $\pm$ 0.0 & 0.0 $\pm$ 0.0 \\
4096 & Depth-OOD & proof-only & 76.1 $\pm$ 0.7 & 26.6 $\pm$ 1.1 & 26.5 $\pm$ 1.1 \\
4096 & Depth-OOD & ProCo & 76.4 $\pm$ 0.6 & 55.1 $\pm$ 0.7 & 55.0 $\pm$ 0.7 \\
\midrule
16384 & ID & answer-only & 94.8 & 0.0 & 0.0 \\
16384 & ID & proof-only & 93.8 & 47.3 & 47.3 \\
16384 & ID & ProCo & 93.7 & 88.5 & 88.4 \\
\cmidrule(lr){2-6}
16384 & Depth-OOD & answer-only & 87.3 & 0.0 & 0.0 \\
16384 & Depth-OOD & proof-only & 85.5 & 37.1 & 37.1 \\
16384 & Depth-OOD & ProCo & 85.0 & 71.2 & 71.2 \\
\midrule
32768 & ID & answer-only & 97.0 & 0.0 & 0.0 \\
32768 & ID & proof-only & 95.6 & 49.0 & 49.0 \\
32768 & ID & ProCo & 96.0 & 93.4 & 93.4 \\
\cmidrule(lr){2-6}
32768 & Depth-OOD & answer-only & 88.7 & 0.0 & 0.0 \\
32768 & Depth-OOD & proof-only & 85.2 & 37.3 & 37.3 \\
32768 & Depth-OOD & ProCo & 87.5 & 77.6 & 77.6 \\
\midrule
112062 & ID & answer-only & 98.5 & 0.0 & 0.0 \\
112062 & ID & proof-only & 98.1 & 51.8 & 51.7 \\
112062 & ID & ProCo & 98.2 & 96.8 & 96.8 \\
\cmidrule(lr){2-6}
112062 & Depth-OOD & answer-only & 93.1 & 0.0 & 0.0 \\
112062 & Depth-OOD & proof-only & 91.3 & 43.6 & 43.6 \\
112062 & Depth-OOD & ProCo & 93.8 & 86.2 & 86.2 \\
\bottomrule
\end{tabular}
\caption{Scaling with larger train budgets.}
\label{tab:scaling-results}
\end{table}

\begin{table}[t]
\centering
\small
\begin{tabular}{rllccc}
\toprule
Train & Eval Set & Variant & Acc. & Verifier-accepted & Joint \\
\midrule
32768 & NatLang Transfer & answer-only & 98.1 & 0.0 & 0.0 \\
32768 & NatLang Transfer & proof-only & 97.4 & 45.1 & 45.1 \\
32768 & NatLang Transfer & \model{} & 97.5 & 93.7 & 93.7 \\
\cmidrule(lr){2-6}
32768 & Birds-Electricity & answer-only & 89.8 & 0.0 & 0.0 \\
32768 & Birds-Electricity & proof-only & 87.2 & 25.2 & 25.2 \\
32768 & Birds-Electricity & \model{} & 91.6 & 90.7 & 90.7 \\
\cmidrule(lr){2-6}
32768 & Depth-3 Transfer & answer-only & 96.5 & 0.0 & 0.0 \\
32768 & Depth-3 Transfer & proof-only & 95.5 & 50.0 & 49.9 \\
32768 & Depth-3 Transfer & \model{} & 95.1 & 92.1 & 92.0 \\
\cmidrule(lr){2-6}
32768 & ID & answer-only & 97.3 & 0.0 & 0.0 \\
32768 & ID & proof-only & 96.4 & 49.3 & 49.2 \\
32768 & ID & \model{} & 96.2 & 93.5 & 93.4 \\
\cmidrule(lr){2-6}
32768 & Depth-OOD & answer-only & 89.3 & 0.0 & 0.0 \\
32768 & Depth-OOD & proof-only & 86.1 & 37.3 & 37.3 \\
32768 & Depth-OOD & \model{} & 87.7 & 77.9 & 77.8 \\
\bottomrule
\end{tabular}
\caption{Transfer across additional ProofWriter test configurations.}
\label{tab:transfer-results}
\end{table}

The transfer results show the same pattern outside the original train/test split. At 32k training questions, \model{} reaches 93.7 joint accuracy on NatLang transfer, 92.0 on depth-3 transfer, and 90.7 on birds-electricity. Raw accuracy stays competitive throughout, and on birds-electricity \model{} is also the strongest raw classifier, reaching 91.6 versus 89.8 for answer-only. Taken together with the full-train depth-OOD result, these numbers suggest that proof/refutation/abstention supervision remains useful under depth, template, and domain shift.
